\section{最大间隔与支持向量机}
\label{sec:09-Convex-Optimization/07-Applications/03}

你有一组带标签 \(y_i\in\{-1,1\}\) 的数据点，想找一个超平面

\[
w^{\trans}x+b=0
\]

按 \(\sign(w^{\trans}x+b)\) 分类。如果数据线性可分，穿过两类之间的分离超平面有无穷多条——稍微转一个角度、平移一截，只要不碰到最近的样本，都算分对。

问题是：挑哪一条？直觉上，你应该挑那条离两类最近的样本都尽可能远的——它左右的活动空间最大，新的测试点稍微偏一点也不至于跨过边界。支持向量机做的就是把这条直觉写成精确的优化问题。

\subsection{固定函数间隔，把几何比例锁住}

点 \(x_i\) 到超平面的有符号几何距离是

\[
\frac{y_i(w^{\trans}x_i+b)}{\norm{ w}_2}.
\]

分子叫函数间隔，分母归一到单位法向量上。这里有个微妙的点：同时把 \((w,b)\) 乘一个正数，超平面纹丝不动（同一个几何位置），但函数间隔跟着变了。如果不把这个自由度锁住，最大化间隔没有确定的意义。

做法是：固定最近样本的函数间隔为 1，

\[
y_i(w^{\trans}x_i+b)\ge 1,
\]

把所有样本都推到支持超平面外面（或至少贴着边界）。在这个约束下，两个支持超平面 \(w^{\trans}x+b=1\) 和 \(w^{\trans}x+b=-1\) 之间的几何距离是 \(2/\norm{ w}_2\)。最大化这个间隔等价于最小化 \(\norm{ w}_2\)，即

\[
\begin{aligned}
\min_{w,b}\quad&\frac12\norm{ w}_2^2\\
\text{s.t.}\quad&
y_i(w^{\trans}x_i+b)\ge 1,\quad i=1,\ldots,n.
\end{aligned}
\]

这是凸二次规划：目标凸、约束仿射。强凸目标保证了 \(w\) 的唯一性，但如果数据的几何分布退化（比如正负样本在某个方向上完全对称），截距 \(b\) 仍可能不唯一。

\subsection{不可分数据：松弛比假装可行诚实}

现实数据通常不能完全线性分开。有些点会掉在间隔内部甚至跨过超平面。强行要求所有点满足硬间隔约束只会导致问题无可行解。引入松弛变量 \(\xi_i\ge0\)：

\[
y_i(w^{\trans}x_i+b)\ge 1-\xi_i,
\]

并惩罚总违规量：

\[
\min_{w,b,\xi}
\frac12\norm{ w}_2^2
+ C\sum_i\xi_i.
\]

消去 \(\xi_i\)（对于每个 \(i\)，最优的 \(\xi_i\) 等于 \(\max\{0,1-y_i(w^{\trans}x_i+b)\}\)）得到 hinge loss 形式：

\[
\min_{w,b}
\frac12\norm{ w}_2^2
+ C\sum_i
\max\{0,1-y_i(w^{\trans}x_i+b)\}.
\]

hinge loss 的行为值得细看。当样本的函数间隔已经大于等于 1 时，loss 为零——分类器对这批点完全满意，不再追加惩罚。间隔不足时线性惩罚，力度由 \(C\) 控制。正是函数间隔超过 1 后那片平坦的零区域，让多数样本在对偶问题中对应的乘子为零——只有贴近边界或跨界的点才需要``用力撑住''超平面。

\(C\) 控制间隔宽度与训练违规的权衡。\(C\) 很大时更在意训练分类正确，可能让超平面对少数异常点过度敏感；\(C\) 很小时更强调大间隔，宁可容忍更多训练违规。\(C\) 需要通过验证来选择，并且受到损失是按样本求和还是按平均的缩放约定影响——改变样本数不改 \(C\) 的含义前，要先确认损失是哪种累法。

\subsection{对偶揭示``支持向量''}

为硬间隔约束引入 Lagrange 乘子 \(\alpha_i\ge0\)。对 \(w,b\) 最小化 Lagrangian，驻点条件给出

\[
w=\sum_i\alpha_i y_i x_i,
\qquad
\sum_i\alpha_i y_i=0.
\]

分类方向是训练样本的加权组合——这已经暗示最后的答案只依赖一部分样本。互补松弛条件给

\[
\alpha_i
\bigl[y_i(w^{\trans}x_i+b)-1\bigr]=0.
\]

若样本严格落在间隔外部——即 \(y_i(w^{\trans}x_i+b) > 1\)——则方括号为正，必有 \(\alpha_i=0\)。只有函数间隔恰好等于 1 的那些点（贴着支持超平面）或软间隔下跨界的点，才可能获得非零乘子。这些点就是支持向量——它们``撑起''了最优超平面的位置和方向。

软间隔的对偶还产生上界

\[
0\le\alpha_i\le C.
\]

这个上界的来历：Lagrangian 中 \(\xi_i\) 的系数为 \(C-\alpha_i\)，同时 \(\xi_i\ge0\) 自带乘子 \(\mu_i\ge0\)。对 \(\xi_i\) 的驻点条件给出 \(C-\alpha_i-\mu_i=0\)，配合 \(\mu_i\ge0\) 就得到 \(\alpha_i\le C\)。\(\alpha_i=C\) 的点是既违规又被限制在惩罚上限的支持向量。

\subsection{核方法仍然靠凸性与半正定}

对偶目标只通过内积 \(x_i^{\trans}x_j\) 使用数据。把内积换成核函数

\[
K(x_i,x_j)=\langle\phi(x_i),\phi(x_j)\rangle
\]

就能在隐式的高维（甚至无穷维）特征空间中寻找线性分隔面——而在原始空间里这个分隔面对应于非线性边界。

核矩阵必须半正定。这不是数学洁癖：只有半正定的核才对应某个特征空间中的内积，才能保持对偶是一个最大化凹二次目标的凸优化问题。用一个不满足 PSD 的相似度函数，二次目标可能不再凹，通用的 SVM 理论收敛结论和求解器假设一起失效。有限精度下出现小负特征值还要结合尺度判断——-1e-10 和 -1e3 含义完全不同。

\subsection{最大间隔不出概率}

SVM 输出的是有符号的 margin \(w^{\trans}x+b\)，一个尺度依赖于 \(\norm{w}\) 的量，并不天然是概率。如果应用需要校准好的概率，可以在独立于训练集的校准数据上做 Platt scaling（用 logistic 函数把 margin 压到 \([0,1]\)）、isotonic regression 等后处理。但这已经是另一个统计模型，SVM 本身的训练过程里没有它。

特征尺度也会改变几何间隔，因为 \(\norm{ w}_2\) 与坐标的单位有关——把``公里''换成``米''，数值间隔变了，各特征的相对重要性在无先验信息时也变了。训练前标准化、明确对截距的处理方式、检查类别权重和报告支持向量占比，都是模型解释的一部分，不是附录里的实现细节。
